\chapter{Cauchy 积分公式}
\label{chap:cauchy-integral-formula-solutions}

\begin{solution}{cauchy-formula-basic-integrals}
Cauchy 公式直接给
\[
 \int_{|z|=R}\frac{e^z}{z-a}\,dz=2\pi i e^a.
\]
高阶导数公式中 \(n=2\) 时
\[
 e^a=\frac{2!}{2\pi i}
 \int_{|z|=R}\frac{e^z}{(z-a)^3}\,dz,
\]
所以第二个积分为 \(\pi i e^a\)。条件 \(|a|<R\) 保证奇点在圆内；若
\(|a|>R\)，两个被积函数都在圆盘内全纯，积分应为零。
\end{solution}

\begin{solution}{cauchy-mean-value}
在
\[
 f(a)=\frac1{2\pi i}\int_{|\zeta-a|=r}
 \frac{f(\zeta)}{\zeta-a}\,d\zeta
\]
中令 \(\zeta=a+re^{it}\)，则
\(d\zeta=ire^{it}dt\)，分母为 \(re^{it}\)，约去后得到
\[
 f(a)=\frac1{2\pi}\int_0^{2\pi}f(a+re^{it})dt.
\]
常数函数显然满足。对 \(f(z)=z^n\)，展开
\[
 (a+re^{it})^n=\sum_{k=0}^n\binom nk a^{n-k}r^ke^{ikt};
\]
积分时 \(k\ge1\) 的项平均为零，只留下 \(a^n=f(a)\)。
\end{solution}

\begin{solution}{cauchy-estimate-coefficients}
系数公式 \(a_n=f^{(n)}(0)/n!\) 与 Cauchy 估计给
\[
 |a_n|=\frac{|f^{(n)}(0)|}{n!}\le\frac M{R^n}.
\]
也可直接使用积分表示
\[
 a_n=\frac1{2\pi i}\int_{|z|=R}\frac{f(z)}{z^{n+1}}dz
\]
和 \(ML\) 估计。必须取 \(R<R_0\)，或另知函数在边界圆的邻域全纯。
\end{solution}

\begin{solution}{isolated-zero-factorization}
Taylor 展开写成
\[
 f(z)=\sum_{n=m}^{\infty}a_n(z-a)^n
 =(z-a)^m\sum_{k=0}^{\infty}a_{m+k}(z-a)^k.
\]
令后一幂级数之和为 \(g\)，则它在同一圆盘全纯，且
\(g(a)=a_m=f^{(m)}(a)/m!\ne0\)。连续性给出一个更小圆盘，使
\(|g(z)-g(a)|<|g(a)|/2\)，故其中 \(g(z)\ne0\)。所以 \(a\) 是
\(f\) 的孤立 \(m\) 阶零点。
\end{solution}

\begin{solution}{entire-polynomial-growth}
取整数 \(n>\alpha\)。在以 \(0\) 为中心、半径 \(R\ge1\) 的圆上，
\[
 M_R\le C(1+R)^\alpha\le C2^\alpha R^\alpha.
\]
Cauchy 估计给
\[
 |f^{(n)}(0)|\le n!C2^\alpha R^{\alpha-n}.
\]
令 \(R\to\infty\)，得到 \(f^{(n)}(0)=0\)。对所有整数
\(n>\alpha\) 都成立，故 \(f\) 的全局 Taylor 级数只含
\(0\le n\le\lfloor\alpha\rfloor\) 的项。全纯函数等于其 Taylor 级数，
所以 \(f\) 是相应次数的多项式。
\end{solution}

\begin{solution}{liouville-real-part}
若 \(\operatorname{Re}f\le M\)，则整函数 \(e^f\) 满足
\[
 |e^{f(z)}|=e^{\operatorname{Re}f(z)}\le e^M.
\]
Liouville 定理使 \(e^f\) 为常数。求导得 \(e^ff'=0\)，而指数从不为零，
故 \(f'=0\)，于是 \(f\) 常值。若实部有下界，对 \(e^{-f}\) 作同样处理。
\end{solution}

\begin{solution}{liouville-bounded-mobius}
像落在闭圆盘 \(\overline{B(w_0,R)}\) 中时，
\(|f(z)|\le |w_0|+R\)，所以 Liouville 定理直接给出常值。

再给出只用最大模原理的证明。取 \(B>0\) 使 \(|f|\le B\)，并定义
\[
 g(z)=
 \begin{cases}
  \dfrac{f(z)-f(0)}z,&z\ne0,\\[4pt]
  f'(0),&z=0.
 \end{cases}
\]
由 \(f\) 在原点的 Taylor 展开，\(g\) 是整函数。给定 \(r>0\)，在
\(|\zeta|=r\) 上有 \(|g(\zeta)|\le 2B/r\)。最大模原理应用于闭圆盘
\(|z|\le r\)，于是
\[
 |g(z)|\le\frac{2B}{r}\qquad(|z|\le r).
\]
固定 \(z\) 并令 \(r\to\infty\)，得到 \(g(z)=0\)。因此
\(f(z)=f(0)\) 对所有 \(z\) 成立。
\end{solution}

\begin{solution}{maximum-modulus-boundary}
\(\overline D\) 是紧集，故连续函数 \(|f|\) 在其上取得最大值。若某个最大点
位于 \(D\)，强最大模原理说明 \(f\) 常值。若 \(f\) 非常数，所有最大点都
属于 \(\overline D\setminus D=\partial D\)，所以
\[
 \max_{\overline D}|f|=\max_{\partial D}|f|.
\]
若 \(f\) 为常数，同一等式仍成立，只是“最大值只能在边界”不再成立。
\end{solution}

\begin{solution}{minimum-modulus-condition}
取单位圆盘上的 \(f(z)=z\)。它非常数，而 \(|f|\) 在内部点 \(0\) 取得
全局最小值 \(0\)。最小模原理的正确说法是：若 \(f\) 全纯、非常数且在一个
邻域不为零，则 \(|f|\) 不能在内部取得正的局部最小值。证明要对 \(1/f\)
应用最大模原理；在零点处 \(1/f\) 没有定义，所以不能省略非零或“正最小值”
条件。
\end{solution}

\begin{solution}{open-mapping-real-valued}
若 \(f\) 非常数，开映射定理说明 \(f(D)\) 是 \(\C\) 中的开集。实轴
\(\R\subset\C\) 不含任何复平面开圆盘，所以非空集合
\(f(D)\subseteq\R\) 不可能在 \(\C\) 中开。矛盾表明 \(f\) 必为常数。
\end{solution}

\begin{solution}{morera-uniform-limit}
紧子集上一致收敛首先给 \(f\) 连续。任取闭含于 \(D\) 的三角形 \(T\)，
其边界紧，故
\[
 \left|\int_{\partial T}(f_n-f)\,dz\right|
 \le \ell(\partial T)\max_{\partial T}|f_n-f|\to0.
\]
每个 \(f_n\) 全纯，所以其边界积分为零，因而
\(\int_{\partial T}f\,dz=0\)。Morera 定理给 \(f\) 全纯。

再取紧集 \(K\Subset D\)，令
\(d=\operatorname{dist}(K,\C\setminus D)>0\)，并在每个 \(z\in K\)
周围取半径 \(d/2\) 的圆。Cauchy 导数公式给
\[
 |f_n^{(k)}(z)-f^{(k)}(z)|
 \le \frac{k!}{(d/2)^k}
 \max_{|\zeta-z|=d/2}|f_n(\zeta)-f(\zeta)|.
\]
这些圆周都落在紧集
\(\{\zeta:\operatorname{dist}(\zeta,K)\le d/2\}\Subset D\)，右端一致
趋零，所以 \(f_n^{(k)}\to f^{(k)}\) 在 \(K\) 上一致。
\end{solution}

\begin{solution}{cauchy-derivative-distance}
令 \(d=\operatorname{dist}(z,\C\setminus D)\)。若 \(d<\infty\)，对每个
\(0<r<d\)，闭圆盘 \(\overline{B(z,r)}\subset D\)，Cauchy 估计给
\[
 |f^{(n)}(z)|\le\frac{n!M}{r^n}.
\]
令 \(r\uparrow d\) 得题中估计。对一般开集，把
\(d(z,\partial D)\) 理解为到补集的距离；当 \(D=\C\) 时补集为空，可约定
\(d=\infty\)。这时对 \(n\ge1\) 令任意圆半径趋于无穷，得到
\(f^{(n)}(z)=0\)，与 Liouville 定理一致。
\end{solution}

\begin{solution}{schwarz-lemma-from-maximum}
定义
\[
 g(z)=
 \begin{cases}f(z)/z,&z\ne0,\\ f'(0),&z=0.\end{cases}
\]
由 \(f\) 在 \(0\) 的 Taylor 展开，\(g\) 在单位圆盘全纯。固定
\(0<r<1\)，最大模原理在 \(|z|\le r\) 上给
\[
 |g(z)|\le\max_{|\zeta|=r}\frac{|f(\zeta)|}{r}\le\frac1r
 \qquad(|z|\le r).
\]
对固定 \(z\) 令 \(r\uparrow1\)，得到 \(|g(z)|\le1\)，所以
\(|f(z)|\le|z|\) 且 \(|f'(0)|=|g(0)|\le1\)。

若某个非零点满足 \(|f(z_0)|=|z_0|\)，或 \(|f'(0)|=1\)，则 \(|g|\)
在内部取得最大值 \(1\)。强最大模原理给 \(g\equiv e^{i\theta}\)，故
\(f(z)=e^{i\theta}z\)。这些旋转也确实处处取等号。
\end{solution}

\begin{solution}{cauchy-estimates-converse}
取实平坦函数
\[
 f(x)=
 \begin{cases}e^{-1/x^2},&x\ne0,\\0,&x=0.\end{cases}
\]
它在 \((-1,1)\) 上 \(C^\infty\)，且 \(f^{(n)}(0)=0\) 对所有 \(n\) 成立。
因此题设中心导数估计对任意 \(M_r>0\) 都成立，但 Taylor 级数恒为零，而
\(f(x)>0\) 对 \(x\ne0\)。中心导数的大小没有控制 Taylor 余项，也没有把
邻域中函数值与这些导数联系起来。

复全纯情形多出的关键条件是 Cauchy 积分表示。它不仅估计中心导数，还通过
\[
 \frac1{\zeta-z}
 =\sum_{n\ge0}\frac{z^n}{\zeta^{n+1}}
\]
的圆周上一致收敛直接证明函数等于所得级数；这一步在一般实
\(C^\infty\) 函数中不存在。
\end{solution}
